\chapter{A Framework for Thinking About Dynamics-Constrained Planning}
\label{ch:framework}

% TODO: Introduce the conceptual lens used to evaluate every contribution
% in the thesis.

\section{Landscape of Approaches}
\label{sec:framework:landscape}

% TODO: Survey planning as search, planning as optimization, and planning
% as inference. Place dynamics-constrained motion generation within this
% taxonomy.

\section{Four Design Axes}
\label{sec:framework:axes}

% TODO: Introduce the four-axis evaluation lens used throughout the thesis.

\subsection{Modeling Burden}
\label{sec:framework:modeling}

% TODO: How much of the world's stochasticity gets baked into the model
% versus offloaded to the planner.

\subsection{Dimensionality and Factorization}
\label{sec:framework:dimensionality}

% TODO: What assumptions and metrics fall out of problem scale, and how
% factorization choices shape tractability.

\subsection{Integration and Pipeline Brittleness}
\label{sec:framework:integration}

% TODO: Error compounding across module boundaries, and the absence of
% mechanisms for continual improvement.

\subsection{Usability}
\label{sec:framework:usability}

% TODO: Scalability, what is worth modeling, generalization, and ease of
% modification when adapting to new tasks.

\section{How These Axes Will Recur}
\label{sec:framework:recurrence}

% TODO: Preview how each subsequent chapter is evaluated against the four
% axes, and why this lens makes the methodological progression legible.
